% 内积空间：内积的定义（公理、例子、复内积的约定）。

\section{内积的定义}
\label{sec:ls-inner-product-axioms}
\phantomsection
\label{sec:ls-inner-product-spaces}
\label{sec:ls-inner-product-definition}

欧氏内积对每个变量都保持线性组合，交换两个变量不改变结果，向量与自身的内积只在零向量处为零。这些性质只涉及线性运算，因此可以用来定义一般空间上的内积。
\begin{definition}[内积空间]
\label{def:functional-inner-product-space}
\label{def:linear-algebra-inner-product}
设$V$是实线性空间。映射$\langle\cdot,\cdot\rangle:V\times V\to\mathbb R$若对任意$\bm x,\bm y,\bm z\in V$与$a,b\in\mathbb R$满足
\begin{enumerate}
 \item 双线性：$\langle a\bm x+b\bm y,\bm z\rangle=a\langle\bm x,\bm z\rangle+b\langle\bm y,\bm z\rangle$，且$\langle\bm z,a\bm x+b\bm y\rangle=a\langle\bm z,\bm x\rangle+b\langle\bm z,\bm y\rangle$；
 \item 对称性：$\langle\bm x,\bm y\rangle=\langle\bm y,\bm x\rangle$；
 \item 正定性：$\langle\bm x,\bm x\rangle\geq0$，且$\langle\bm x,\bm x\rangle=0$当且仅当$\bm x=\bm0$，
\end{enumerate}
就称为$V$上的\term{内积}（inner product）；配备内积的线性空间$(V,\langle\cdot,\cdot\rangle)$称为\term{内积空间}（inner-product space）。
\end{definition}

\paragraph{内积的例子}
先检验欧氏内积满足这三条公理。设$\bm x,\bm y,\bm z\in\mathbb R^d$，$a,b\in\mathbb R$。对称性成立，因为每一项$x_jy_j=y_jx_j$。对第一个变量的线性来自有限和的运算法则：
\[
 \sum_j(ax_j+bz_j)y_j=a\sum_jx_jy_j+b\sum_jz_jy_j,
\]
再由对称性，对第二个变量同样成立。正定性成立，因为$\langle\bm x,\bm x\rangle=\sum_jx_j^2\geq0$，而非负数之和为零时每一项都为零。

同形矩阵上的Frobenius内积也满足三条公理。它对两个变量分别线性；交换两个矩阵不改变各项乘积，因而对称；$\langle\bm A,\bm A\rangle_{\mathrm F}=\sum_{i,j}a_{i,j}^2$非负，并且只在$\bm A=\bm0$时为零。这正是把$mn$个元素看作坐标时欧氏内积的三条性质。所以下面从公理推出的一般结论对它同样成立。

标准内积之外，一般空间上还可以有各种内积。下面在$P_1$上规定一个只看两点取值的内积，后面的算例会反复用到它。
\begin{example}[只看两点取值的内积]
\label{ex:ls-two-point-inner-product}
在$P_1$上规定
\[
 \langle f,g\rangle=f(0)g(0)+f(1)g(1),
\]
即把两个函数在$t=0,1$处的取值配对相乘再相加。函数值$f(0)$、$f(1)$都随$f$线性变化，所以右端对$f$与对$g$分别线性；交换$f,g$不改变右端，对称性也成立。

正定性要用到$P_1$的特点。由
\[
 \langle f,f\rangle=f(0)^2+f(1)^2\geq0,
\]
若它为零，则$f(0)=f(1)=0$；对$f=a+bt$，这就是$a=0$、$a+b=0$，故$f=0$。一次多项式由两个点的取值确定，这一步才能成立。

因此这一规则是$P_1$上的内积。它与系数向量$(a,b)^\top$的欧氏内积并不相同：例如$\langle1,t\rangle=1\cdot0+1\cdot1=1$，而$1$与$t$的系数向量$(1,0)^\top$、$(0,1)^\top$的欧氏内积为零。同一对函数是否正交，要看采用哪种内积。
\end{example}

\paragraph{复内积的约定}
上面的定义针对实空间。在复空间中，若两个变量都按复数线性，并且满足实数式的对称性，那么$\langle\mathrm i\bm x,\mathrm i\bm x\rangle=-\langle\bm x,\bm x\rangle$，无法对非零向量保持正定。因此复内积采用另一种约定：对第一个变量共轭线性，对第二个变量线性，并满足$\langle\bm x,\bm y\rangle=\overline{\langle\bm y,\bm x\rangle}$。标准例子是$\langle\bm x,\bm y\rangle=\bm x^*\bm y$，其中$\bm x^*$表示共轭转置；仍要求$\langle\bm x,\bm x\rangle>0$对非零$\bm x$成立。

后面的Cauchy--Schwarz不等式在复内积下同样成立。若$\bm y\ne\bm0$，在$\lVert\bm x-\alpha\bm y\rVert^2\geq0$中取$\alpha=\langle\bm y,\bm x\rangle/\lVert\bm y\rVert^2$，展开可得
\[
 |\langle\bm x,\bm y\rangle|^2\leq\lVert\bm x\rVert^2\lVert\bm y\rVert^2 .
\]
因此它的平方根仍给出范数。

但复内积的值可能不是实数，下一节用余弦定义夹角的公式不能直接照搬。下文的证明继续使用实内积，复情形需要把转置相应改为共轭转置。
